\begin{proofbox}
\textbf{证明前统一记号：}
设 $F_1(-c,0)$、$F_2(c,0)$，$P(x_P,y_P)$ 在椭圆上且不为长轴顶点；
$r_1=PF_1$、$r_2=PF_2$、$\theta=\angle F_1PF_2$。
必有
\[
r_1+r_2=2a,\qquad F_1F_2=2c,\qquad
b^2=a^2-c^2,\qquad 0<\theta<\pi.
\]

\textbf{方法一：代数法——余弦定理与面积公式（完整通法）}

\textbf{第 1 步：} 在 $\triangle F_1PF_2$ 中使用余弦定理：
\[
4c^2=r_1^2+r_2^2-2r_1r_2\cos\theta.
\]
因为 $r_1^2+r_2^2=(r_1+r_2)^2-2r_1r_2$，所以
\[
\begin{aligned}
4c^2
 &=4a^2-2r_1r_2(1+\cos\theta)\\
 &=4a^2-4r_1r_2\cos^2\frac\theta2.
\end{aligned}
\]
移项后得到第一个核心恒等式：
\[
\boxed{r_1r_2\cos^2\frac\theta2=a^2-c^2=b^2}.\tag{1}
\]
\textbf{第 2 步：} 同一个三角形有两种面积写法：
\[
S=\frac12r_1r_2\sin\theta
 =\frac12(2c)|y_P|=c|y_P|.
\]
由 $\sin\theta=2\sin(\theta/2)\cos(\theta/2)$，
\[
r_1r_2\sin\frac\theta2\cos\frac\theta2=c|y_P|.
\]
两边除以式 (1)（非退化时 $\cos(\theta/2)>0$），得
\[
\boxed{\tan\frac\theta2=\frac{c|y_P|}{b^2}}.\tag{2}
\]
\textbf{第 3 步：} 对式 (1) 使用均值不等式：
\[
r_1r_2\le\left(\frac{r_1+r_2}{2}\right)^2=a^2.
\]
因此
\[
\cos^2\frac\theta2=\frac{b^2}{r_1r_2}\ge\frac{b^2}{a^2},
\]
\[
\sin^2\frac\theta2\le1-\frac{b^2}{a^2}
=\frac{c^2}{a^2}=e^2.
\]
因为 $\sin(\theta/2)>0$，
\[
\boxed{\sin\frac\theta2\le e}.
\]
等号成立当且仅当 $r_1=r_2=a$；这意味着 $P$ 在
$F_1F_2$ 的垂直平分线上，故 $x_P=0$；再代入椭圆方程得
$P=(0,\pm b)$。故
\[
\boxed{\theta_{\max}=2\arcsin e}.
\]

\textbf{方法二：几何法——内切圆、角平分线与切线长（不设坐标求焦半径）}

给焦点三角形作内切圆，设圆心为 $I$，半径为 $\rho$。
从 $I$ 向边 $PF_1$ 作垂线，垂足为 $T$。
$PI$ 平分 $\angle F_1PF_2$，因而在直角三角形 $PIT$ 中
\[
\tan\frac\theta2=\frac{IT}{PT}=\frac\rho{PT}.\tag{3}
\]
三角形的半周长为
\[
s=\frac{r_1+r_2+2c}{2}=a+c.
\]
利用同一点引圆的两条切线段相等，可知\textbf{顶点 $P$ 的切线段长}
\[
PT=\frac{r_1+r_2-2c}{2}=a-c.\tag{4}
\]
另一方面，三角形面积既是底乘高的一半，也是内切圆半径乘半周长：
\[
S=\frac12(2c)h=ch=\rho s,\qquad h=|y_P|.
\]
因此 $\rho=ch/(a+c)$，代入式 (3)、(4)：
\[
\tan\frac\theta2
=\frac{ch/(a+c)}{a-c}
=\frac{ch}{(a+c)(a-c)}
=\boxed{\frac{ch}{b^2}}.
\]
这是从\textbf{内切圆 + 等切线长 + 面积}直接走到半角正切式的纯几何路线。
最后，由椭圆几何性质 $0<h\le b$，
且 $\tan t$ 在 $0<t<90^\circ$ 递增，
可知 $\theta$ 在 $h=b$ 的两个短轴顶点取得最大值。
此时 $PF_1=PF_2=a$，等腰三角形的顶角平分线构成直角三角形，
$\sin(\theta_{\max}/2)=c/a=e$。

\textbf{方法三：直观证明——高度控制视角（先图形，后严格收口）}

固定焦点间距 $2c$，并把 $P$ 看作移动的观察者。
椭圆定义固定了 $PF_1+PF_2=2a$；点到焦点连线的高度记为 $h$。
在长轴顶点附近 $h\approx0$，两个焦点在观察者眼中几乎重合，
$\theta\approx0$；在短轴顶点 $h=b$，视角最张开。

这个图形观察可用方法二得到的\textbf{严格单调关系}封口：
\[
\theta=2\arctan\frac{ch}{b^2},\qquad 0<h\le b.
\]
因 $c/b^2>0$ 且 $\arctan$ 严格递增，
\[
\theta\le2\arctan\frac cb.
\]
在 $P=(0,\pm b)$ 时取等号。再由
$a^2=b^2+c^2$ 和直角三角形，
\[
\sin\left(\arctan\frac cb\right)=\frac c{\sqrt{b^2+c^2}}=\frac ca=e.
\]
于是 $\sin(\theta/2)\le e$。这也解释了：
\textbf{最大张角由最大垂直高度控制，而不是由某条焦半径单独控制。}

\textbf{方法四：高考最快法——三行推出最大角}

若只需 $\sin(\theta/2)\le e$ 与等号条件，不必求坐标，也不必求面积：
\[
\begin{aligned}
4c^2
 &= (r_1+r_2)^2-4r_1r_2\cos^2\frac\theta2\\
 &=4a^2-4r_1r_2\cos^2\frac\theta2;
\end{aligned}
\]
\[
\boxed{\cos^2\frac\theta2=\frac{b^2}{r_1r_2}
\ge\frac{b^2}{a^2}}
\quad\Longrightarrow\quad
\boxed{\sin\frac\theta2\le\frac ca=e}.
\]
等号 $\Longleftrightarrow r_1=r_2=a
\Longleftrightarrow P=(0,\pm b)$。

\textbf{补充：解析几何坐标交叉检验（发现正切倒写错误的可靠办法）}

由焦半径公式
\[
r_1=a+\frac ca x_P,\qquad
r_2=a-\frac ca x_P,
\]
得到 $r_1r_2=a^2-c^2x_P^2/a^2$，且
\[
\tan^2\frac\theta2
=\frac{r_1r_2-b^2}{b^2}
=\frac{c^2(a^2-x_P^2)}{a^2b^2}.
\]
由椭圆方程 $a^2-x_P^2=a^2y_P^2/b^2$，
\[
\tan^2\frac\theta2=\frac{c^2y_P^2}{b^4}.
\]
半角在 $(0,90^\circ)$，正切为正，故再次得到
$\tan(\theta/2)=c|y_P|/b^2$。
\end{proofbox}
